\documentclass[11pt]{article}
\usepackage[utf8]{inputenc}
\usepackage[T1]{fontenc}
\usepackage[a4paper,margin=1in]{geometry}
\usepackage{parskip}
\usepackage{hyperref}
\hypersetup{hidelinks}
\pagestyle{empty}

\begin{document}

\noindent\textbf{\large Cover letter --- submission to \emph{Universe} (MDPI)}

\medskip
\noindent
\begin{tabular}{@{}ll@{}}
\textbf{Manuscript:} & A Coulomb--Newton Cosmology: Based on a \\
                     & Proton--Electron Model (RealQM\,/\,RealNucleus) \\
\textbf{Author:}     & Claes Johnson, KTH Royal Institute of Technology, Stockholm, Sweden \\
\textbf{Type:}       & Article (theoretical / exploratory) \\
\textbf{Section:}    & Cosmology \\
\end{tabular}

\bigskip
\hrule
\bigskip

\noindent Dear Editors,

\medskip
I am pleased to submit the manuscript \emph{A Coulomb--Newton Cosmology: Based on a Proton--Electron Model
(RealQM\,/\,RealNucleus)} for consideration as an Article in \emph{Universe}.

The paper sets out an exploratory, non-standard cosmological model built from only two well-understood forces ---
electron--proton Coulomb interaction and Newtonian gravitation --- with no Big Bang singularity, no inflation, no
dark-matter particle, and no cosmological constant. The organising idea is that Coulomb and Newton are the
\emph{same} field equation (Poisson's) with opposite sign and a different strength: a two-valued Poisson structure
with two sources (charge and mass/energy) and two constants, one per force. Coulomb, acting on charge, builds a
small-scale proton--electron world of nuclei and atoms (the RealQM\,/\,RealNucleus program, no strong force);
Newton, acting on the energy of that matter, gathers it at the large scale, its positive and negative curvature
sectors separating into the cosmic web and a repulsive (dark-energy-like) component. The model reproduces a helium
mass fraction near $Y\approx0.25$ as a consistency check. The two forces are treated in parallel; the paper is
careful to leave open whether they ultimately reduce to one.

The manuscript is deliberately framed as a proposal, not as established cosmology. It includes an explicit
Status note, a section relating the model to $\Lambda$CDM and to observations (stating plainly what it does and
does not yet address --- the redshift--distance relation, the CMB, precise large-scale-structure statistics, and
light elements beyond helium), and a candid list of load-bearing open problems, chiefly what sets the nuclear
energy scale. What the model offers in exchange for its incompleteness is economy: fewer free primitives than the
standard account.

I believe the work fits the scope of \emph{Universe}, which welcomes original and exploratory theoretical
contributions to cosmology and foundational physics, including approaches outside the mainstream that are argued
carefully and honestly. The paper is self-contained and does not presuppose familiarity with the author's earlier
RealQM and RealNucleus work, which is cited for context.

The manuscript is original, has not been published elsewhere, and is not under consideration by another journal.
The author has no competing interests to declare. There are no funding sources to acknowledge and no data sets
associated with the work (it is analytical/theoretical; supporting interactive simulations are freely available
online and are cited).

Thank you for considering the manuscript. I would be happy to respond to any questions.

\bigskip
\noindent Yours sincerely,

\medskip
\noindent Claes Johnson \\
Professor emeritus of Applied Mathematics, KTH Royal Institute of Technology \\
\texttt{claesjohnson@gmail.com}

\end{document}
